\setAuthor{Taavet Kalda}
\setRound{piirkonnavoor}
\setYear{2019}
\setNumber{G 8}
\setDifficulty{7}
\setTopic{Dünaamika}

\prob{Granaat}
Granaat visatakse algkiirusega $v$ õhku nurga $\alpha$ all. Hetkel $t_1$ plahvatab granaat $N\gg1$ killuks, kusjuures granaadi massikeskme süsteemis eemalduvad killud kiirusega $u$ ühtlase nurkjaotusega. Mis ajahetkel $t_2$ jõuab esimene kild maapinnani ning mis ajahetkel $t_3$ viimane? Raskuskiirendus on $g$. 

\hint
Peale plahvatust võime ette kujutada, et granaat lendab edasi, kuid seda ümbritseb kildudest koosnev sfäär raadiusega $u(t-t_1)$. Teame, et esimene ja viimane kild vastavad sfääril vastavalt kõige alumise ja ülemise punktiga.\solu
Granaadi vertikaalsuunaline algkiirus on $v_0=v\sin\alpha$.

Kuna massikeskme süsteemis kaugenevad killud võrdse kiirusega $u$ keskpunktist, siis võime ette kujutada, et ajahetkel $t$ peale granaadi lõhkemist lendab ta edasi, kuid seda ümbritseb kildudest koosnev sfäär raadiusega $ut$. Paneme tähele, et sfääri keskpunkt liigub nii, nagu oleks liikunud granaat, kui see ei oleks lõhkenud.

Kui esimene kild jõuab maapinaanni, on sfääri raadius $u(t_2-t_1)$, ehk ajahetkel $t_2$ on sfääri keskpunkti kõrguseks $u(t_2-t_1)$. Sfääri keskpunkti liikumisvõrrandist on selle lennukõrgus ajahetkel $t_2$ antud kui $h_2=v_0t_2-gt_2^2/2$, ehk
$$
v_0t_2-\frac{gt_2^2}{2}=u(t_2-t_1).
$$
Selle võrrandi positiivne lahend on
\[
t_2=\frac{(v_0-u)+\sqrt{(v_0-u)^2+2gt_1u}}{g}.
\]

Kui viimane kild jõuab maapinaanni, on sfääri (mis nüüdseks on juba mõtteline, sest ülejäänud selle sfääri pinnal olevad killud lebavad maas) raadius $u(t_3-t_1)$. Seega, ajahetkel $t_3$ on meie mõttelise sfääri keskpunkt kõrgusel $h_3=-u(t_3-t_1)$. See annab võrrandiks
$$
v_0t_3-\frac{gt_3^2}{2}=-u(t_3-t_1).
$$
Pidades silmas, et $t_3>t_2$, saame selle lahendiks
\[
t_3=\frac{v_0+u+\sqrt{(v_0+u)^2-2gt_1u}}{g}.
\]\probend